% \section{Testing Datasets}
% The datasets described below are used for the experiments outlined in the following sections (see Sections \ref{sec:processing_app}, and \ref{sec:tracking_app})
% \begin{itemize}
%     \item \textbf{Generic Sine Waves with Noise}: To validate the correctness of our processing 
%     algorithms as well as their speed, we generate binary files of synthetic radar data consisting of 
%     a complex sine wave with additive white Gaussian noise. This allows us to test the software under 
%     controlled conditions where we can precompute the expected output, and to evaluate the performance 
%     of the software in terms of latency and throughput without the confounding factors of real-world 
%     data. (INCLUDE FIGURE HERE illustrating the signal)  
%     \item \textbf{AWR2243 Recorded Data}: To validate the operation of our pipeline with realistic 
%     data speeds, timing uncertainties, and electrical noise, we record data from our AWR2243 radar 
%     system in TEST\_SOURCE mode. In test source mode, two targets can be simulated on the chip with 
%     position velocity and bounds. The receivers record what the reflected signals from these targets 
%     would look like. While the data recorded is not fully representative of real-world data, it allows 
%     us to test the software with realistic data rates and noise characteristics. 
% \end{itemize}

\section{Testing Datasets}
The datasets described below are used for the experiments outlined in the previous sections 
(\ref{sec:processing_app} and \ref{sec:tracking_app}).

The first dataset is used 
to validate both the correctness and the speed of our processing algorithms during development. 
By generating binary data streams containing complex sine waves with Gaussian noise using bash scripts, 
we can test the radar software under controlled conditions where the expected output can be precomputed 
in advance. This makes it possible to verify algorithmic correctness while also evaluating latency and 
throughput without the additional complexity introduced by real-world phenomena such as multipath 
propagation, data corruption, and non-gaussian noise. An interactive variation of this dataset is 
also used for debugging, allowing the user to adjust parameters such as SNR, range, and bearing/elevation
ratios in real. Figure \ref{fig:synthetic-signal} illustrates an example of a synthetic signal generated 
for testing.

\begin{figure}
    \incfig[0.6]{dataaset1}
    \caption[Example of a synthetic signal generated for testing.]{An 
    example of a synthetic signal generated for testing. The top plot shows the real components of the 
    samples of many received chirps (stacked vertically). Fast time increases to the right and the height 
    of a sample within a chirp is its voltage. In the bottom plot, the horizontal axis is now range, 
    successive chirps are stacked vertically and amplitude is depicted with colour, where yellow is the highest 
    amplitude. 2nd plot is produced by applying an \gls{fft} to each synthetic chirp from the first plot. 
    The transform allow us to visually validate that the frequency to range conversion is correct, 
    and the preprocessing stage is behaving as expected.}
    \label{fig:synthetic-signal}
\end{figure}

The second dataset consists of recorded data from the AWR2243 radar system operating in \verb|TEST_SOURCE| 
mode. This dataset is used to validate the behaviour of the processing pipeline under more realistic 
data rates, timing uncertainties, and electrical noise conditions. In \verb|TEST_SOURCE| mode, the chip can 
simulate two targets with configurable position, velocity, and bounds, and the receivers capture the 
signals that would be reflected from these targets. Although this recorded data is still not fully 
representative of real-world radar returns, it provides a more realistic approximation of practical 
operating conditions than the synthetic sine wave dataset. 